% !PRINTED 439
\weberpage{439}
\begin{center}
\S.\,138.

\gesperrt{Vorzeichenbestimmung. Quadratische Reste.}
\end{center}

Wir haben im vorigen Paragraphen zwei Formeln abgeleitet, in denen noch Vorzeichen zu bestimmen waren, nämlich
\[
2^{\frac{n-1}{2}} \overset{\nu}{\Pi} \cos\frac{2\nu m\pi}{n} = \pm 1, \tag{1}
\]
\[
2^{\frac{n-1}{2}} \overset{\nu}{\Pi} \sin\frac{2\nu m\pi}{n} = \pm \sqrt{n}, \tag{2}
\]
worin $m$ relativ prim zu der ungeraden Zahl $n$ ist und $\nu$ die Reihe der Zahlen
\[
1,\ 2,\ \cdots\ \frac{n-1}{2} \tag{$\nu$}
\]
durchläuft. Wir wollen das Zeichen $\nu$ für die Zahlen dieser Reihe hier festhalten.

Jede beliebige ganze Zahl giebt bei der Theilung durch $n$ als Rest eine der Zahlen $\nu$ oder $0$ oder $n-\nu$. Statt $n-\nu$ können wir, wenn wir auch negative Reste zulassen, $-\nu$ wählen, und wenn wir also von den durch $n$ theilbaren Zahlen absehen, so bleibt eine der Zahlen $\pm\nu$ als Rest. Wir nennen diese den \gesperrt{absolut kleinsten Rest} (zum Unterschied von dem kleinsten Rest im gewöhnlichen Sinne, der aus den Zahlen $1, 2, \ldots n-1$ genommen ist).

Es sei nun also $m$ eine zu $n$ theilerfremde Zahl; wir betrachten die Reihe der Zahlen
\[
m,\ 2m,\ 3m \cdots \frac{n-1}{2}\, m \tag{$m\nu$}
\]
und bilden zu jeder den absolut kleinsten Rest $\varrho$:
\[
\pm \nu_1,\ \pm \nu_2,\ \pm \nu_3 \cdots \pm \nu_{\frac{1}{2}(n-1)}; \tag{$\varrho$}
\]
unter diesen Zahlen $\varrho$ kommen nicht zwei gleiche und auch nicht zwei entgegengesetzte vor; denn wenn die Summe oder die Differenz zweier Zahlen $(\varrho)$, also $\varrho+\varrho'$ oder $\varrho-\varrho'$ gleich Null wäre, so müsste für zwei verschiedene Zahlen $\nu, \nu'$ aus $(\nu)$
\[
m(\nu \pm \nu')
\]
durch $n$ theilbar sein, also müsste auch $\nu \pm \nu'$ durch $n$ theilbar sein, was unmöglich ist, da jede der Zahlen $\nu, \nu'$ kleiner als $n:2$ ist,
%CONT>>
